% ============================================================================
% CamBus - Manual de Usuario y Documentación Técnica
% Centro de Distribución Mabe SLP
% ============================================================================
% Documento LaTeX para Overleaf
% ============================================================================

\documentclass[12pt,a4paper]{article}

% ============================================================================
% PAQUETES
% ============================================================================
\usepackage[utf8]{inputenc}
\usepackage[spanish]{babel}
\usepackage{graphicx}
\usepackage{xcolor}
\usepackage{listings}
\usepackage{hyperref}
\usepackage{geometry}
\usepackage{fancyhdr}
\usepackage{titlesec}
\usepackage{tocloft}
\usepackage{booktabs}
\usepackage{array}
\usepackage{tabularx}
\usepackage{longtable}
\usepackage{enumitem}
\usepackage{tikz}
\usepackage{float}
\usepackage{caption}
\usepackage{subcaption}
\usepackage{awesomebox}
\usepackage{fontawesome5}

% ============================================================================
% CONFIGURACIÓN DE PÁGINA
% ============================================================================
\geometry{
    left=2.5cm,
    right=2.5cm,
    top=3cm,
    bottom=3cm
}

% Colores corporativos
\definecolor{mabe-blue}{RGB}{31, 119, 180}
\definecolor{mabe-dark}{RGB}{38, 39, 48}
\definecolor{code-bg}{RGB}{240, 242, 246}
\definecolor{success}{RGB}{40, 167, 69}
\definecolor{warning}{RGB}{255, 193, 7}
\definecolor{danger}{RGB}{220, 53, 69}

% Configuración de hyperref
\hypersetup{
    colorlinks=true,
    linkcolor=mabe-blue,
    filecolor=mabe-blue,
    urlcolor=mabe-blue,
    citecolor=mabe-blue,
    pdftitle={CamBus - Manual de Usuario},
    pdfauthor={CamBus Team},
    pdfsubject={Documentación del Sistema CamBus}
}

% Configuración de listings para código
\lstset{
    backgroundcolor=\color{code-bg},
    basicstyle=\ttfamily\small,
    breaklines=true,
    captionpos=b,
    commentstyle=\color{gray},
    frame=single,
    keepspaces=true,
    keywordstyle=\color{mabe-blue}\bfseries,
    numbers=left,
    numbersep=5pt,
    numberstyle=\tiny\color{gray},
    rulecolor=\color{black},
    showspaces=false,
    showstringspaces=false,
    showtabs=false,
    stringstyle=\color{success},
    tabsize=4,
    xleftmargin=1cm
}

% Estilos de código específicos
\lstdefinestyle{bash}{
    language=bash,
    morekeywords={sudo, apt, pacman, pip, docker, streamlit}
}

\lstdefinestyle{yaml}{
    basicstyle=\ttfamily\small,
    morecomment=[l]{\#}
}

% Configuración de encabezados
\pagestyle{fancy}
\fancyhf{}
\fancyhead[L]{\leftmark}
\fancyhead[R]{\includegraphics[height=0.8cm]{example-image}}
\fancyfoot[C]{\thepage}
\renewcommand{\headrulewidth}{0.4pt}
\renewcommand{\footrulewidth}{0.4pt}

% Formato de secciones
\titleformat{\section}
    {\normalfont\Large\bfseries\color{mabe-blue}}
    {\thesection}{1em}{}
\titleformat{\subsection}
    {\normalfont\large\bfseries\color{mabe-dark}}
    {\thesubsection}{1em}{}

% ============================================================================
% DOCUMENTO
% ============================================================================

\begin{document}

% ============================================================================
% PORTADA
% ============================================================================
\begin{titlepage}
    \centering
    \vspace*{2cm}
    
    {\Huge\bfseries\color{mabe-blue} CamBus}\\[0.5cm]
    {\LARGE Control de Acceso Vehicular}\\[1cm]
    
    \rule{\textwidth}{0.5pt}\\[0.5cm]
    {\Large\bfseries Manual de Usuario y Documentación Técnica}\\[0.3cm]
    \rule{\textwidth}{0.5pt}\\[2cm]
    
    {\large Centro de Distribución Mabe SLP}\\[3cm]
    
    \begin{tikzpicture}
        \node[draw=mabe-blue, line width=2pt, rounded corners=10pt, inner sep=20pt, fill=code-bg] {
            \Large\faIcon{truck} \quad Sistema de Gestión de Andenes
        };
    \end{tikzpicture}
    
    \vfill
    
    {\large Versión 1.0.0}\\[0.5cm]
    {\large Febrero 2026}
    
\end{titlepage}

% ============================================================================
% ÍNDICE
% ============================================================================
\tableofcontents
\newpage

% ============================================================================
% 1. INTRODUCCIÓN
% ============================================================================
\section{Introducción}

\subsection{Descripción del Sistema}

\textbf{CamBus} es una aplicación web desarrollada con Streamlit que permite la gestión integral de entrada y salida de vehículos en el Centro de Distribución Mabe SLP. El sistema está diseñado para controlar y monitorear 100 andenes de carga/descarga en tiempo real.

\subsection{Características Principales}

\begin{itemize}[leftmargin=2cm]
    \item[\faIcon{tachometer-alt}] \textbf{Dashboard en Tiempo Real}: Visualización del estado de 100 andenes con actualización automática cada 5 segundos
    \item[\faIcon{clipboard-list}] \textbf{Registro Manual}: Entrada y salida de vehículos de forma manual
    \item[\faIcon{chart-bar}] \textbf{Reportes y Análisis}: Flujo por hora, eficiencia de andenes, tiempos de permanencia
    \item[\faIcon{cogs}] \textbf{Administración}: Gestión de usuarios, andenes, cámaras y backups
\end{itemize}

\subsection{Requisitos del Sistema}

\begin{table}[H]
    \centering
    \begin{tabular}{|l|l|}
        \hline
        \textbf{Componente} & \textbf{Versión Mínima} \\
        \hline
        Python & 3.9+ \\
        PostgreSQL & 12+ \\
        Docker (opcional) & 20.10+ \\
        \hline
    \end{tabular}
    \caption{Requisitos mínimos del sistema}
\end{table}

% ============================================================================
% 2. INSTALACIÓN
% ============================================================================
\section{Instalación}

\subsection{Instalación en Linux - Distribución Debian/Ubuntu}

\subsubsection{Paso 1: Actualizar el Sistema}

\begin{lstlisting}[style=bash]
sudo apt update && sudo apt upgrade -y
\end{lstlisting}

\subsubsection{Paso 2: Instalar Dependencias del Sistema}

\begin{lstlisting}[style=bash]
# Instalar Python y pip
sudo apt install -y python3 python3-pip python3-venv

# Instalar PostgreSQL
sudo apt install -y postgresql postgresql-contrib

# Instalar dependencias adicionales
sudo apt install -y build-essential libpq-dev git curl
\end{lstlisting}

\subsubsection{Paso 3: Configurar PostgreSQL}

\begin{lstlisting}[style=bash]
# Iniciar el servicio de PostgreSQL
sudo systemctl start postgresql
sudo systemctl enable postgresql

# Crear base de datos y usuario
sudo -u postgres psql -c "CREATE DATABASE cambus_db;"
sudo -u postgres psql -c "ALTER USER postgres PASSWORD 'tu_password';"

# Cargar el esquema de la base de datos
sudo -u postgres psql -d cambus_db -f cambus.sql
\end{lstlisting}

\subsubsection{Paso 4: Configurar la Aplicación}

\begin{lstlisting}[style=bash]
# Clonar o copiar el proyecto
cd /ruta/al/proyecto/cambus

# Crear entorno virtual
python3 -m venv venv
source venv/bin/activate

# Instalar dependencias de Python
pip install -r requirements.txt

# Configurar variables de entorno
cp .env.example .env
nano .env  # Editar con tus credenciales
\end{lstlisting}

\subsubsection{Paso 5: Ejecutar la Aplicación}

\begin{lstlisting}[style=bash]
# Activar entorno virtual (si no está activo)
source venv/bin/activate

# Ejecutar la aplicación
streamlit run app.py
\end{lstlisting}

La aplicación estará disponible en: \url{http://localhost:8501}

\subsection{Instalación en Linux - Distribución Arch Linux}

\subsubsection{Paso 1: Actualizar el Sistema}

\begin{lstlisting}[style=bash]
sudo pacman -Syu
\end{lstlisting}

\subsubsection{Paso 2: Instalar Dependencias del Sistema}

\begin{lstlisting}[style=bash]
# Instalar Python y pip
sudo pacman -S python python-pip

# Instalar PostgreSQL
sudo pacman -S postgresql

# Instalar dependencias adicionales
sudo pacman -S base-devel git curl
\end{lstlisting}

\subsubsection{Paso 3: Configurar PostgreSQL}

\begin{lstlisting}[style=bash]
# Inicializar el cluster de PostgreSQL
sudo -u postgres initdb -D /var/lib/postgres/data

# Iniciar el servicio
sudo systemctl start postgresql
sudo systemctl enable postgresql

# Crear base de datos
sudo -u postgres createdb cambus_db

# Establecer contraseña
sudo -u postgres psql -c "ALTER USER postgres PASSWORD 'tu_password';"

# Cargar esquema
sudo -u postgres psql -d cambus_db -f cambus.sql
\end{lstlisting}

\subsubsection{Paso 4: Configurar la Aplicación}

\begin{lstlisting}[style=bash]
# Navegar al proyecto
cd /ruta/al/proyecto/cambus

# Crear entorno virtual
python -m venv venv
source venv/bin/activate

# Instalar dependencias
pip install -r requirements.txt

# Configurar variables de entorno
cp .env.example .env
vim .env  # Editar con tus credenciales
\end{lstlisting}

\subsubsection{Paso 5: Ejecutar la Aplicación}

\begin{lstlisting}[style=bash]
source venv/bin/activate
streamlit run app.py
\end{lstlisting}

\subsection{Instalación en Windows}

\subsubsection{Paso 1: Instalar Python}

\begin{enumerate}
    \item Descargar Python 3.11+ desde \url{https://www.python.org/downloads/}
    \item Ejecutar el instalador
    \item \textbf{Importante}: Marcar la opción ``Add Python to PATH''
    \item Completar la instalación
\end{enumerate}

\subsubsection{Paso 2: Instalar PostgreSQL}

\begin{enumerate}
    \item Descargar PostgreSQL desde \url{https://www.postgresql.org/download/windows/}
    \item Ejecutar el instalador
    \item Configurar la contraseña del usuario \texttt{postgres}
    \item Recordar el puerto (por defecto: 5432)
\end{enumerate}

\subsubsection{Paso 3: Crear la Base de Datos}

Abrir \textbf{pgAdmin} o usar la línea de comandos:

\begin{lstlisting}[style=bash]
# Abrir PowerShell o CMD
psql -U postgres

# En el prompt de PostgreSQL
CREATE DATABASE cambus_db;
\q

# Cargar el esquema
psql -U postgres -d cambus_db -f cambus.sql
\end{lstlisting}

\subsubsection{Paso 4: Configurar la Aplicación}

\begin{lstlisting}[style=bash]
# Abrir PowerShell y navegar al proyecto
cd C:\ruta\al\proyecto\cambus

# Crear entorno virtual
python -m venv venv
.\venv\Scripts\activate

# Instalar dependencias
pip install -r requirements.txt

# Copiar y editar configuración
copy .env.example .env
notepad .env
\end{lstlisting}

\subsubsection{Paso 5: Ejecutar la Aplicación}

\begin{lstlisting}[style=bash]
# Activar entorno virtual
.\venv\Scripts\activate

# Ejecutar la aplicación
streamlit run app.py
\end{lstlisting}

\subsection{Instalación con Docker (Todas las Plataformas)}

Para una instalación simplificada usando contenedores:

\begin{lstlisting}[style=bash]
# Clonar el proyecto
git clone <repositorio> cambus
cd cambus

# Configurar contraseña de base de datos
export DB_PASSWORD=tu_password_seguro

# Construir y ejecutar
docker-compose up --build -d

# Verificar que los contenedores están corriendo
docker-compose ps
\end{lstlisting}

La aplicación estará disponible en:
\begin{itemize}
    \item \textbf{CamBus}: \url{http://localhost:8501}
    \item \textbf{Adminer} (opcional): \url{http://localhost:8080}
\end{itemize}

% ============================================================================
% 3. ESTRUCTURA DEL PROYECTO
% ============================================================================
\section{Estructura del Proyecto}

\subsection{Árbol de Directorios}

\begin{lstlisting}[basicstyle=\ttfamily\footnotesize]
cambus/
|-- app.py                 # Aplicacion principal (punto de entrada)
|-- config.yaml            # Configuracion de la aplicacion
|-- requirements.txt       # Dependencias de Python
|-- cambus.sql             # Esquema de base de datos PostgreSQL
|-- Dockerfile             # Imagen Docker
|-- docker-compose.yml     # Orquestacion de contenedores
|-- .env.example           # Plantilla de variables de entorno
|-- utils/
|   |-- __init__.py
|   |-- db.py              # Conexion y funciones de base de datos
|   |-- auth.py            # Autenticacion y manejo de sesiones
|-- components/
|   |-- __init__.py
|   |-- sidebar.py         # Barra lateral de navegacion
|-- pages/
|   |-- __init__.py
|   |-- dashboard.py       # Dashboard principal con KPIs
|   |-- registro.py        # Registro manual de vehiculos
|   |-- reportes.py        # Reportes y analisis
|   |-- admin.py           # Administracion del sistema
\end{lstlisting}

\subsection{Descripción de Componentes}

\begin{table}[H]
    \centering
    \begin{tabularx}{\textwidth}{|l|X|}
        \hline
        \textbf{Archivo/Directorio} & \textbf{Descripción} \\
        \hline
        \texttt{app.py} & Punto de entrada principal. Configura la página, estilos CSS, y enruta a las diferentes páginas según la selección del usuario. \\
        \hline
        \texttt{config.yaml} & Archivo de configuración con parámetros de base de datos, dashboard, estados de andenes y roles de usuario. \\
        \hline
        \texttt{utils/db.py} & Módulo de conexión a PostgreSQL. Contiene todas las funciones de consulta y modificación de datos. \\
        \hline
        \texttt{utils/auth.py} & Sistema de autenticación. Maneja login, sesiones, hash de contraseñas (bcrypt) y control de permisos por rol. \\
        \hline
        \texttt{components/sidebar.py} & Componente de navegación lateral. Muestra menú según permisos del usuario autenticado. \\
        \hline
        \texttt{pages/dashboard.py} & Dashboard principal con KPIs, mapa de andenes en grid 10x10, y lista de vehículos activos. \\
        \hline
        \texttt{pages/registro.py} & Formularios para registrar entrada y salida manual de vehículos. Validación de placas mexicanas. \\
        \hline
        \texttt{pages/reportes.py} & Generación de reportes: flujo horario, eficiencia de andenes, tiempos de permanencia. Exportación a CSV/Excel. \\
        \hline
        \texttt{pages/admin.py} & Panel administrativo para gestión de usuarios, andenes, cámaras y backups de base de datos. \\
        \hline
    \end{tabularx}
    \caption{Descripción de los componentes del sistema}
\end{table}

\subsection{Base de Datos}

El sistema utiliza PostgreSQL con las siguientes tablas principales:

\begin{table}[H]
    \centering
    \begin{tabularx}{\textwidth}{|l|X|}
        \hline
        \textbf{Tabla} & \textbf{Descripción} \\
        \hline
        \texttt{turnos} & Turnos de trabajo (Matutino, Vespertino, Nocturno) \\
        \hline
        \texttt{usuarios} & Usuarios del sistema con roles y credenciales \\
        \hline
        \texttt{andenes} & Catálogo de 100 andenes con estado y zona \\
        \hline
        \texttt{camaras} & Cámaras de vigilancia asociadas a andenes \\
        \hline
        \texttt{registros\_vehiculos} & Registro de entrada/salida (tabla particionada por año) \\
        \hline
        \texttt{estancias\_vehiculos} & Historial de estancias con duración calculada \\
        \hline
        \texttt{productos} & Catálogo de productos para carga/descarga \\
        \hline
        \texttt{carga\_descarga\_productos} & Movimientos de productos por estancia \\
        \hline
        \texttt{bitacora\_acciones} & Log de auditoría de acciones de usuario \\
        \hline
    \end{tabularx}
    \caption{Tablas principales de la base de datos}
\end{table}

% ============================================================================
% 4. FUNCIONALIDADES
% ============================================================================
\section{Funcionalidades del Sistema}

\subsection{Dashboard en Tiempo Real}

El dashboard principal proporciona una vista general del centro de distribución:

\begin{itemize}
    \item \textbf{KPIs Principales}:
    \begin{itemize}
        \item Total de andenes (100)
        \item Andenes ocupados / libres
        \item Vehículos ingresados hoy
        \item Tiempo promedio de estancia
    \end{itemize}
    
    \item \textbf{Mapa de Andenes}: Grid visual de 10x10 mostrando el estado de cada andén:
    \begin{itemize}
        \item[\textcolor{success}{\faIcon{circle}}] Verde: Libre
        \item[\textcolor{danger}{\faIcon{circle}}] Rojo: Ocupado
        \item[\textcolor{warning}{\faIcon{circle}}] Amarillo: En mantenimiento
        \item[\faIcon{circle}] Negro: Bloqueado
    \end{itemize}
    
    \item \textbf{Auto-refresh}: Actualización automática cada 5 segundos
    
    \item \textbf{Lista de Vehículos Activos}: Tabla con vehículos actualmente en andenes
\end{itemize}

\subsection{Registro Manual de Vehículos}

Permite el registro manual de entrada y salida de vehículos:

\subsubsection{Registro de Entrada}

\begin{enumerate}
    \item Ingresar la placa del vehículo (formato mexicano)
    \item Seleccionar el andén de destino (solo andenes libres)
    \item Seleccionar el tipo de vehículo:
    \begin{itemize}
        \item TORTON
        \item RABÓN
        \item FULL
        \item CAMIONETA
    \end{itemize}
    \item Opcionalmente ingresar nombre del transportista
    \item Confirmar el registro
\end{enumerate}

\subsubsection{Registro de Salida}

\begin{enumerate}
    \item Seleccionar el vehículo activo a dar salida
    \item Agregar observaciones (opcional)
    \item Confirmar la salida
\end{enumerate}

\subsection{Reportes y Análisis}

El sistema genera los siguientes tipos de reportes:

\begin{table}[H]
    \centering
    \begin{tabularx}{\textwidth}{|l|X|}
        \hline
        \textbf{Reporte} & \textbf{Descripción} \\
        \hline
        Flujo por Hora & Análisis del flujo de vehículos por hora del día. Gráficos de líneas mostrando entradas y salidas. \\
        \hline
        Eficiencia de Andenes & Métricas por andén: operaciones totales, tiempo promedio, mínimo, máximo y mediana. \\
        \hline
        Tiempos de Permanencia & Distribución de tiempos de estancia. Histogramas y estadísticas descriptivas. \\
        \hline
        Productos Más Movidos & Ranking de productos con mayor movimiento de carga/descarga. \\
        \hline
    \end{tabularx}
    \caption{Tipos de reportes disponibles}
\end{table}

\textbf{Exportación}: Todos los reportes pueden exportarse en formato CSV o Excel.

\subsection{Administración del Sistema}

El panel de administración (solo para rol ADMIN) permite:

\subsubsection{Gestión de Usuarios}

\begin{itemize}
    \item Ver listado de usuarios con estado y último acceso
    \item Crear nuevos usuarios con roles asignados
    \item Activar/desactivar usuarios
    \item Cambiar contraseñas
\end{itemize}

\subsubsection{Gestión de Andenes}

\begin{itemize}
    \item Ver estado de todos los andenes
    \item Cambiar estado manualmente (LIBRE, OCUPADO, MANTENIMIENTO, BLOQUEADO)
    \item Filtrar por zona
\end{itemize}

\subsubsection{Gestión de Cámaras}

\begin{itemize}
    \item Listado de cámaras instaladas
    \item Agregar nuevas cámaras
    \item Modificar configuración (IP, modelo, estado)
    \item Eliminar cámaras
\end{itemize}

\subsubsection{Backups}

\begin{itemize}
    \item Generar backup de la base de datos
    \item Descargar archivos de respaldo
\end{itemize}

% ============================================================================
% 5. ROLES Y PERMISOS
% ============================================================================
\section{Roles y Permisos}

El sistema implementa control de acceso basado en roles (RBAC):

\begin{table}[H]
    \centering
    \begin{tabular}{|l|c|c|c|c|}
        \hline
        \textbf{Rol} & \textbf{Dashboard} & \textbf{Registro} & \textbf{Reportes} & \textbf{Admin} \\
        \hline
        ADMIN & \checkmark & \checkmark & \checkmark & \checkmark \\
        \hline
        SUPERVISOR & \checkmark & \checkmark & \checkmark & $\times$ \\
        \hline
        OPERADOR & \checkmark (solo lectura) & $\times$ & $\times$ & $\times$ \\
        \hline
    \end{tabular}
    \caption{Matriz de permisos por rol}
\end{table}

\subsection{Seguridad}

\begin{itemize}
    \item \textbf{Contraseñas}: Almacenadas con hash bcrypt (12 rounds)
    \item \textbf{Sesiones}: Manejadas con \texttt{st.session\_state}
    \item \textbf{Intentos fallidos}: Máximo 5 intentos antes de bloqueo temporal
    \item \textbf{Validación}: Control de permisos en cada operación
\end{itemize}

% ============================================================================
% 6. CONFIGURACIÓN
% ============================================================================
\section{Configuración}

\subsection{Archivo config.yaml}

\begin{lstlisting}[style=yaml]
# Configuracion de Base de Datos
database:
  host: "localhost"
  port: 5432
  name: "cambus_db"
  user: "postgres"
  password: ""  # Usar variable de entorno

# Configuracion del Dashboard
dashboard:
  refresh_interval: 5
  total_andenes: 100
  grid_columns: 10
  grid_rows: 10

# Estados de Andenes
andenes:
  estados:
    LIBRE: "verde"
    OCUPADO: "rojo"
    MANTENIMIENTO: "amarillo"
    BLOQUEADO: "negro"
\end{lstlisting}

\subsection{Variables de Entorno}

\begin{table}[H]
    \centering
    \begin{tabular}{|l|l|l|}
        \hline
        \textbf{Variable} & \textbf{Descripción} & \textbf{Default} \\
        \hline
        \texttt{DB\_PASSWORD} & Contraseña de PostgreSQL & - \\
        \hline
        \texttt{DB\_HOST} & Host de PostgreSQL & localhost \\
        \hline
        \texttt{DB\_PORT} & Puerto de PostgreSQL & 5432 \\
        \hline
        \texttt{DB\_NAME} & Nombre de la base de datos & cambus\_db \\
        \hline
        \texttt{DB\_USER} & Usuario de PostgreSQL & postgres \\
        \hline
    \end{tabular}
    \caption{Variables de entorno del sistema}
\end{table}

% ============================================================================
% 7. COMANDOS ÚTILES
% ============================================================================
\section{Comandos Útiles}

\subsection{Desarrollo}

\begin{lstlisting}[style=bash]
# Ejecutar con auto-reload
streamlit run app.py --server.runOnSave true

# Ejecutar en puerto especifico
streamlit run app.py --server.port 8080

# Modo debug
streamlit run app.py --logger.level debug
\end{lstlisting}

\subsection{Base de Datos}

\begin{lstlisting}[style=bash]
# Generar backup
pg_dump -h localhost -U postgres cambus_db > backup.sql

# Restaurar backup
psql -h localhost -U postgres cambus_db < backup.sql

# Conectar a la base de datos
psql -h localhost -U postgres -d cambus_db
\end{lstlisting}

\subsection{Docker}

\begin{lstlisting}[style=bash]
# Construir imagen
docker-compose build

# Iniciar servicios
docker-compose up -d

# Ver logs
docker-compose logs -f cambus-app

# Detener servicios
docker-compose down

# Reiniciar aplicacion
docker-compose restart cambus-app

# Incluir Adminer
docker-compose --profile tools up -d
\end{lstlisting}

% ============================================================================
% 8. SOLUCIÓN DE PROBLEMAS
% ============================================================================
\section{Solución de Problemas}

\subsection{Errores Comunes}

\begin{table}[H]
    \centering
    \begin{tabularx}{\textwidth}{|l|X|}
        \hline
        \textbf{Error} & \textbf{Solución} \\
        \hline
        ``Connection refused'' a PostgreSQL & Verificar que el servicio PostgreSQL está corriendo: \texttt{sudo systemctl status postgresql} \\
        \hline
        ``ModuleNotFoundError'' & Activar el entorno virtual: \texttt{source venv/bin/activate} \\
        \hline
        ``Permission denied'' en puerto 8501 & Usar un puerto diferente o ejecutar con permisos elevados \\
        \hline
        Error de autenticación BD & Verificar credenciales en \texttt{.env} y permisos del usuario PostgreSQL \\
        \hline
    \end{tabularx}
    \caption{Errores comunes y soluciones}
\end{table}

\subsection{Verificación de Servicios}

\begin{lstlisting}[style=bash]
# Verificar PostgreSQL
sudo systemctl status postgresql
psql -U postgres -c "SELECT 1;"

# Verificar puerto disponible
lsof -i :8501

# Verificar conexion a base de datos
python -c "from utils.db import test_connection; print(test_connection())"
\end{lstlisting}

% ============================================================================
% APÉNDICE
% ============================================================================
\appendix

\section{Dependencias de Python}

\begin{lstlisting}[basicstyle=\ttfamily\small]
# Framework Web
streamlit>=1.28.0

# Base de Datos
psycopg2-binary>=2.9.9
sqlalchemy>=2.0.0

# Procesamiento de Datos
pandas>=2.0.0
numpy>=1.24.0

# Visualizacion
plotly>=5.17.0

# Autenticacion y Seguridad
bcrypt>=4.0.0

# Configuracion
python-dotenv>=1.0.0
pyyaml>=6.0.0

# Exportacion de Reportes
openpyxl>=3.1.0
xlsxwriter>=3.1.0
\end{lstlisting}

% ============================================================================
% FIN DEL DOCUMENTO
% ============================================================================

\vfill

\begin{center}
    \rule{0.5\textwidth}{0.5pt}\\[0.5cm]
    {\large\textbf{CamBus - Control de Acceso Vehicular}}\\
    {\small Centro de Distribución Mabe SLP}\\[0.3cm]
    {\small © 2024-2026 CamBus Team}
\end{center}

\end{document}
